\documentclass[aspectratio=169,11pt]{beamer}

\usepackage{fontspec}
\usepackage[ngerman,provide=*]{babel}
\usepackage{csquotes}

\usetheme{default}
\usecolortheme{dove}
\setbeamertemplate{navigation symbols}{}
\setbeamerfont{title}{size=\Large,series=\bfseries}
\setbeamerfont{frametitle}{size=\large,series=\bfseries}
\setbeamercolor{frametitle}{fg=black}
\setbeamercolor{title}{fg=black}
\setbeamercolor{structure}{fg=black!75}
\setbeamercolor{normal text}{fg=black!90}
\setbeamertemplate{itemize item}{\textbullet}
\setbeamertemplate{itemize subitem}{\textendash}

\setbeamertemplate{footline}{%
  \leavevmode%
  \hbox{%
    \begin{beamercolorbox}[wd=.55\paperwidth,ht=2.25ex,dp=1ex,leftskip=1em]{author in head/foot}%
      \footnotesize\color{black!60}MTA -- Sitzung 3: Vom Verfahren zur Interpretation%
    \end{beamercolorbox}%
    \begin{beamercolorbox}[wd=.45\paperwidth,ht=2.25ex,dp=1ex,rightskip=1em,leftskip=2ex]{title in head/foot}%
      \hfill\footnotesize\color{black!60}\insertframenumber\,/\,\inserttotalframenumber%
    \end{beamercolorbox}}%
}

\usepackage{booktabs}
\usepackage{array}
\usepackage{tabularx}
\usepackage{amsmath}
\usepackage{tikz}

\title{Vom Verfahren zur Interpretation}
\subtitle{Sitzung 3 -- Vorverarbeitung, Metadaten\\und ein kritischer Blick auf Topic-Modelle}
\author{Prof.~Dr.~Christian Papilloud}
\institute{Institut für Soziologie\\Martin-Luther-Universität Halle-Wittenberg}
\date{Master Soziologie -- 2026}

\begin{document}

% ===============================================================
\begin{frame}[plain]
  \titlepage
\end{frame}

% ===============================================================
\begin{frame}{Wo wir stehen}

  In den ersten beiden Sitzungen haben wir Topic-Modelle
  \emph{methodologisch} positioniert und \emph{technisch}
  verstanden.

  \vfill
  Heute schließen wir den theoretischen Teil ab mit dem, was
  zwischen \emph{Forschungsfrage} und \emph{Algorithmus} liegt --
  und was \emph{nach} dem Algorithmus folgt:

  \vfill
  \begin{enumerate}
    \item[\textbf{1.}] \emph{Vor} dem Algorithmus:
          Was ist ein Dokument? Wie reinigt man eine Textquelle?
          Was ist TF-IDF?
    \item[\textbf{2.}] Wie kann man \emph{Metadaten} in eine
          Topic-Modell-Analyse einbinden? Drei Strategien.
    \item[\textbf{3.}] \emph{Nach} dem Algorithmus -- der kritische
          Blick: Zwei verbreitete Erwartungen, die Topic-Modelle
          \emph{nicht} erfüllen.
    \item[\textbf{4.}] Was sich für unsere Forschungspraxis ändert.
          Übergang zur Praxis mit MTA.
  \end{enumerate}

\end{frame}

% ===============================================================
\section{1. Was \emph{vor} dem Algorithmus geschieht}

\begin{frame}{Die unsichtbare Hauptarbeit}

  Eine Topic-Modell-Analyse besteht aus vier Phasen
  (Papilloud \& Hinneburg 2018, Kap.\ 2.4):

  \vfill
  \begin{enumerate}
    \item \textbf{Datenvorverarbeitung} \hfill -- der heutige Schwerpunkt
    \item \textbf{Topic-Modell-Inferenz} \hfill -- letzte Sitzung
    \item \textbf{Evaluation}
    \item \textbf{Interpretation}
  \end{enumerate}

  \vfill
  \begin{block}{Pointe}
    Die Vorverarbeitung ist der Schritt, der die meiste Zeit kostet --
    und der die größten methodologischen Konsequenzen hat. Wer schlecht
    vorverarbeitet, bekommt schlechte topics. Es gibt keinen Algorithmus,
    der diese Phase ersetzt.
  \end{block}

\end{frame}

% ===============================================================
\begin{frame}{Frage 1: Was ist ein \emph{Dokument}?}

  Es klingt trivial, ist es aber nicht.

  \vfill
  \begin{itemize}
    \item Ein Interview = \emph{ein} Dokument? Oder eine
          \emph{Aussage} = ein Dokument?
    \vfill
    \item Eine ganze Tageszeitung = ein Dokument? Oder ein
          \emph{Artikel}? Oder ein \emph{Absatz}?
    \vfill
    \item Ein Buch = ein Dokument? Oder ein \emph{Kapitel}?
          Oder eine \emph{Seite}?
  \end{itemize}

  \vfill
  Diese Wahl ist \emph{nicht} technisch. Sie ist eine \emph{theoretische
  Wahl}: Auf welcher Ebene sehen Sie die Strukturen, die Sie suchen?

  \vfill
  Beispiel: Wenn Sie tageszeitungsweise gruppieren, finden Sie topics,
  die ganze \emph{Ausgaben} dominieren. Wenn Sie artikelweise gruppieren,
  finden Sie topics, die einzelne \emph{Themen} ausmachen. Wenn Sie
  absatzweise gruppieren, finden Sie \emph{rhetorische Muster}. Drei
  Analysen, drei \emph{verschiedene} Befunde.

\end{frame}

% ===============================================================
\begin{frame}{Frage 2: Welche Wörter zählen?}

  Nicht alle Wörter sind gleich nützlich für die Themenfindung.
  Zwei Gruppen werden meistens entfernt:

  \vfill
  \begin{enumerate}
    \item \textbf{Wenig aussagekräftige Wörter} (Stoppwörter):
          Artikel, Pronomen, Präpositionen, Hilfsverben.
          \\[2pt]
          \emph{ein, mein, er, sie, und, aber, ist, hat, durch, mit \dots}
          \\[2pt]
          Sehr häufig, semantisch arm.
    \vfill
    \item \textbf{Sehr seltene Wörter}: Wörter, die nur in einem oder
          zwei Dokumenten vorkommen.
          \\[2pt]
          Sie können keinem topic stabil zugewiesen werden und
          \emph{stören} die Inferenz statistisch.
  \end{enumerate}

  \vfill
  \footnotesize\color{black!60}
  In MTA: Sie laden eine Stoppwörterliste hoch
  (\texttt{stopwords\_de.txt}, \texttt{stopwords\_en.txt}, \dots).
  Bei der Datenvorbereitung können Sie zusätzliche Wörter manuell
  ausschließen.

\end{frame}

% ===============================================================
\begin{frame}{TF-IDF -- eine zentrale Gewichtung}

  TF-IDF (Term Frequency -- Inverse Document Frequency) gewichtet
  Wörter doppelt:

  \vfill
  \begin{itemize}
    \item \textbf{TF} -- Häufigkeit eines Wortes \emph{im Dokument}.
          Je öfter, desto wichtiger für dieses Dokument.
    \vfill
    \item \textbf{IDF} -- inverse Häufigkeit \emph{über die Dokumente}.
          \[ \mathrm{idf}_i = \log_{10}\!\left(\frac{N}{n_i}\right) \]
          mit $N$ = Gesamtzahl der Dokumente, $n_i$ = Anzahl der
          Dokumente, die das Wort $i$ enthalten.
          \\[2pt]
          Wörter, die in \emph{allen} Dokumenten vorkommen, bekommen
          einen niedrigen IDF-Wert.
  \end{itemize}

  \vfill
  \textbf{Beispiel:} im Reuters-Korpus (806.791 Texte) haben die
  Wörter \emph{Versuch} und \emph{Versicherung} ähnliche
  Gesamthäufigkeiten ($\sim$10\,400). Aber \emph{Versuch} kommt in
  8.760 Dokumenten vor, \emph{Versicherung} nur in 3.997.
  \emph{Versicherung} ist also stärker konzentriert -- höherer
  IDF-Wert -- für die Topic-Analyse aussagekräftiger.

\end{frame}

% ===============================================================
\begin{frame}{Lemmatisierung, NER, Phrasen -- optionale Schritte}

  Drei sprachliche Operationen, die Topic-Modelle verbessern können:

  \vfill
  \begin{itemize}
    \item \textbf{Lemmatisierung.} Reduziert Wörter auf ihre Grundform:
          \emph{lief}, \emph{laufen}, \emph{gelaufen} $\rightarrow$ \emph{laufen}.
          \\[2pt]
          Vermeidet, dass dieselbe semantische Einheit über mehrere
          topics verteilt wird.
    \vfill
    \item \textbf{Named Entity Recognition (NER).} Erkennt
          Eigennamen und mehrwortige Ausdrücke:
          \\[2pt]
          \emph{Martin Luther University Halle-Wittenberg}
          $\rightarrow$ ein Token statt fünf.
    \vfill
    \item \textbf{Wortphrasen.} Häufige Wortgruppen werden zu
          \emph{einem} Token zusammengefasst.
          \\[2pt]
          \emph{battery life} $\rightarrow$ \emph{battery\_life}.
  \end{itemize}

  \vfill
  \textbf{In MTA:} Stoppwörter und TF-IDF sind eingebaut.
  Lemmatisierung, NER und Phrasen \emph{nicht} -- sie können vorher
  durch externe Werkzeuge (TreeTagger, spaCy) angewendet werden,
  wenn nötig.

\end{frame}

% ===============================================================
\begin{frame}{Eine wichtige Empfehlung: pro Sprache, pro Korpus}

  \begin{block}{Aus dem Buch (Papilloud \& Hinneburg 2018, S.~8 der Doc)}
    \enquote{Wenn Sie Dateien in verschiedenen Sprachen haben, legen
    Sie sie \emph{nicht} in dasselbe Verzeichnis. Erstellen Sie
    so viele Verzeichnisse wie nötig, z.B.\ eines für Ihre
    französischen Dateien, eines für Ihre deutschen, eines für
    Ihre englischen.}
  \end{block}

  \vfill
  Diese Trennung gilt aber nicht nur für Sprachen. Sie gilt auch für:
  \begin{itemize}
    \item Verschiedene \emph{Autoren} (wenn der Stil unterschiedlich ist),
    \vfill
    \item Verschiedene \emph{Genres} (Interviews vs.\ Briefe vs.\ Berichte),
    \vfill
    \item Verschiedene \emph{Zeitperioden} (wenn das Vokabular sich
          stark gewandelt hat).
  \end{itemize}

  \vfill
  Jede Mischung führt zu \emph{vermischten} topics. Trennen, dann
  vergleichen.

\end{frame}

% ===============================================================
\section{2. Metadaten und Topic-Modelle}

\begin{frame}{Was sind Metadaten?}

  \emph{Meta\-daten} = Informationen über die Dokumente, die nicht
  im Text enthalten sind:
  \vfill
  \begin{itemize}
    \item Zeit der Veröffentlichung,
    \item Autorin oder Autor,
    \item Geschlecht, Alter, Bildungsstand der Befragten,
    \item Quelle (Zeitung, Plattform, Region),
    \item Zitate, Genre, Sprache \ldots
  \end{itemize}

  \vfill
  Diese Informationen sind für die soziologische Interpretation
  \emph{wertvoll}. Wie binden wir sie in eine Topic-Modell-Analyse ein?

  \vfill
  Es gibt drei Strategien.

\end{frame}

% ===============================================================
\begin{frame}{Strategie 1 -- \emph{Spaltung}}

  \textbf{Idee:} Man teilt das Korpus nach der Metadaten-Variable
  in Teilkorpora und führt für jedes eine separate Topic-Analyse durch.

  \vfill
  \emph{Beispiel:} Korpus aus 2.000 Zeitungsartikeln zu Migration,
  2015--2020. Wir wollen die \emph{zeitliche Entwicklung} sehen.

  $\rightarrow$ Wir machen sechs separate Analysen, eine pro Jahr,
  und vergleichen die topics.

  \vfill
  \textbf{Vorteil:} einfach umzusetzen, sauber interpretierbar pro
  Teilgruppe.

  \vfill
  \textbf{Nachteil:} jede Teilgruppe muss groß genug bleiben.
  Bei wenigen Dokumenten pro Jahr ist die Analyse instabil.
  Außerdem: man verliert die \emph{gemeinsamen} topics über die
  Gruppen hinweg.

  \vfill
  \footnotesize\color{black!60}
  \emph{In MTA} entspricht das: man legt mehrere Verzeichnisse an
  und analysiert sie nacheinander.

\end{frame}

% ===============================================================
\begin{frame}{Strategie 2 -- \emph{Hinzufügung}}

  \textbf{Idee:} Die Metadaten werden \emph{nach} der Topic-Inferenz
  in die Analyse eingebunden. Man rechnet \emph{ein} Topic-Modell auf
  dem gesamten Korpus und betrachtet dann, wie sich die
  Topic-Verteilungen \emph{nach Gruppen unterscheiden}.

  \vfill
  \emph{Beispiel:} Wir berechnen 8 topics auf allen Migrations-Artikeln.
  Dann fragen wir: kommt Topic 3 (Wirtschaftliche Aspekte) bei der FAZ
  häufiger vor als bei der taz?

  \vfill
  \textbf{Vorteil:} flexibel, statistisch gut kontrollierbar (man
  kann statistische Tests anwenden, wie in MTA mit Welch + BH).

  \vfill
  \textbf{Nachteil:} die topics selbst sind \emph{nicht} durch die
  Metadaten beeinflusst. Wenn die Metadaten die topics
  \emph{strukturieren}, geht das verloren.

  \vfill
  \footnotesize\color{black!60}
  \emph{In MTA} entspricht das: \enquote{Group Comparison}-Menü
  (mit Welch-t-Test und Benjamini-Hochberg-Korrektur).

\end{frame}

% ===============================================================
\begin{frame}{Strategie 3 -- \emph{Lernergänzung}}

  \textbf{Idee:} Die Metadaten werden \emph{während} der Topic-Inferenz
  einbezogen. Das Topic-Modell selbst wird so erweitert, dass es die
  Metadaten als Teil seines Modells nutzt.

  \vfill
  Beispiele aus der Literatur:
  \begin{itemize}
    \item \emph{Dynamic Topic Models} (Blei \& Lafferty 2006):
          topics, die sich über die Zeit verändern dürfen.
    \item \emph{Author-Topic Models} (Rosen-Zvi et al.\ 2004):
          Autoren haben eigene Topic-Profile.
    \item \emph{Structural Topic Model (STM)} (Roberts et al.\ 2014):
          Metadaten beeinflussen sowohl Topic-Verteilung als auch
          Topic-Inhalt.
  \end{itemize}

  \vfill
  \textbf{Vorteil:} elegant; topics werden \emph{strukturell} mit
  Metadaten verbunden.

  \vfill
  \textbf{Nachteil:} mathematisch komplex; spezialisierte Implementierungen
  nötig. \emph{In MTA nicht implementiert.}

\end{frame}

% ===============================================================
\begin{frame}{Welche Strategie wann?}

  \footnotesize
  \begin{tabularx}{\textwidth}{@{}p{0.20\textwidth}XXX@{}}
    \toprule
    & \textbf{Spaltung} & \textbf{Hinzufügung} & \textbf{Lernergänzung} \\
    \midrule
    \textbf{Voraussetzung}
      & große Teilgruppen
      & gemeinsames Korpus
      & spezialisiertes Modell \\[6pt]
    \textbf{Topic-Inhalt}
      & kann pro Gruppe variieren
      & gemeinsam für alle Gruppen
      & strukturell beeinflusst \\[6pt]
    \textbf{Interpretation}
      & gruppenweise
      & vergleichend
      & vereinheitlicht \\[6pt]
    \textbf{Statistische Tests}
      & schwer möglich
      & \textbf{ja} (Welch, Mann-Whitney, BH)
      & im Modell integriert \\[6pt]
    \textbf{Verfügbar in MTA}
      & ja
      & ja
      & nein \\
    \bottomrule
  \end{tabularx}

  \vfill
  \normalsize
  \footnotesize\color{black!60}
  Empfehlung des Buches (Papilloud \& Hinneburg 2018, S.~35):
  Strategie 2 (Hinzufügung) ist meistens vorzuziehen, wenn die
  Metadaten nicht zwingend strukturell modelliert werden müssen.

\end{frame}

% ===============================================================
\section{3. Der kritische Blick}

\begin{frame}{Topic-Modelle haben Grenzen}

  Wir kommen zum Kapitel 8 des Buches -- \emph{Rück- und Ausblick}.
  Zwei verbreitete Erwartungen werden dort explizit zurückgewiesen.
  Beide werden Sie in der Literatur und in der Praxis immer wieder
  hören. Beide sind \emph{falsch}.

  \vfill
  \begin{itemize}
    \item \textbf{Erwartung 1:} Topic-Modelle überwinden die
          Unterscheidung qualitativ / quantitativ.
    \vfill
    \item \textbf{Erwartung 2:} Topic-Modelle reduzieren die
          Interpretationsarbeit \emph{auf ein Minimum}.
  \end{itemize}

  \vfill
  Wir nehmen beide Erwartungen kurz auseinander.

\end{frame}

% ===============================================================
\begin{frame}{Erwartung 1 -- \enquote{Quali/Quanti überwinden}}

  Das Argument lautet meistens: Da Topic-Modelle Statistik und
  Mathematik einsetzen, lösen sie das alte Paar
  \emph{qualitativ vs.\ quantitativ} auf.

  \vfill
  \textbf{Warum das nicht stimmt} (Buch, S.~152):
  \begin{itemize}
    \item Topic-Modelle setzen die Arbeit an \emph{Texten} voraus --
          also am \emph{Sinn} von Begriffen.
    \vfill
    \item Diese Daten lassen sich nicht auf eine einzige Bedeutung
          reduzieren. Sie können nicht rein quantitativ interpretiert
          werden.
    \vfill
    \item Topic-Modelle reduzieren die Komplexität von Textquellen
          \emph{nicht} -- im Gegenteil, sie laden dazu ein,
          \emph{Multiperspektivität} zu erkennen.
  \end{itemize}

  \vfill
  \begin{block}{Pointe}
    Topic-Modelle können quantitative Verfahren \emph{begleiten}
    (Mixed-Methods-Forschung). Aber sie sind nicht \emph{a priori}
    \enquote{besser} oder \enquote{neutraler}, weil sie mit Zahlen
    arbeiten.
  \end{block}

\end{frame}

% ===============================================================
\begin{frame}{Erwartung 2 -- \enquote{Interpretation auf ein Minimum reduzieren}}

  Das Argument: Da der Algorithmus die Kategorien selbst findet,
  müssen wir nur noch die Ergebnisse \enquote{lesen}.

  \vfill
  \textbf{Warum das nicht stimmt} (Buch, S.~152--153):
  \begin{itemize}
    \item Die Ergebnisse eines Topic-Modells hängen \emph{stark} ab von:
          \begin{itemize}
            \item der Forschungsfrage (was wir suchen);
            \item den theoretischen Rahmenannahmen;
            \item den Vorverarbeitungsentscheidungen;
            \item der Wahl von $K$;
            \item der Wahl von LDA oder NMF;
            \item der Stoppwörterliste;
            \item \dots
          \end{itemize}
    \vfill
    \item Die topics sind \emph{Hypothesen}, keine Befunde.
    \vfill
    \item Die Bedeutung der topics muss durch Lektüre der besten
          Dokumente, durch Vergleich mit Theorie, durch Diskussion
          mit Kollegen \emph{erarbeitet} werden.
  \end{itemize}

\end{frame}

% ===============================================================
\begin{frame}{Eine Pflichtlektüre: Schmidt 2012}

  \begin{block}{Schmidt 2012 -- die Kritik}
    Schmidt zeigt anhand historischer und geographischer Daten, wie
    Topic-Modelle zu \emph{falschen Ergebnissen} führen können,
    wenn sie \enquote{blind} angewendet werden.
  \end{block}

  \vfill
  Seine Forderung: Topic-Modelle in den Geistes- und Sozialwissenschaften
  müssen \textbf{mit Vorsicht und kritischer Reflexivität} angewendet
  werden.

  \vfill
  Was heißt das konkret?
  \begin{enumerate}
    \item Topics nicht für bare Münze nehmen -- immer mit den
          besten Dokumenten konfrontieren.
    \item Mehrere Modelle laufen lassen (verschiedene $K$, LDA und
          NMF) und Konsistenz prüfen.
    \item Ergebnisse mit unabhängigen Quellen vergleichen.
    \item Die theoretische Rahmung explizit machen.
  \end{enumerate}

\end{frame}

% ===============================================================
\begin{frame}{Wo bleibt der/die Forschende?}

  \begin{block}{Frage vom Ende der ersten Sitzung}
    Eine Sozialwissenschaftlerin, die mit MTA arbeitet, codiert nicht
    selbst, sondern sie liest die Codierung einer Maschine. Verschiebt
    das die Subjektivität der Interpretation -- oder verlagert es sie
    nur an einen anderen Ort?
  \end{block}

  \vfill
  Die Antwort lautet jetzt: \emph{ja, es verlagert sie}. Aber:

  \vfill
  \begin{itemize}
    \item Die Verlagerung ist \emph{sichtbarer}: alle Vorverarbeitungs- und
          Modellierungsentscheidungen sind dokumentierbar.
    \vfill
    \item Sie ist \emph{prüfbar}: jede andere Forscherin kann die
          Analyse mit denselben Parametern wiederholen.
    \vfill
    \item Sie ist \emph{kommunizierbar}: man kann mit Kollegen über
          Parameter streiten, was bei rein interpretativen Verfahren
          schwerer ist.
  \end{itemize}

  \vfill
  Die Subjektivität verschwindet nicht. Sie wird transparenter.

\end{frame}

% ===============================================================
\section{4. Was sich für unsere Praxis ändert}

\begin{frame}{Eine andere Art zu lesen}

  Das vielleicht wichtigste, was Topic-Modelle in die qualitative
  Forschung einbringen, ist eine \emph{neue Lesepraxis}:

  \vfill
  \begin{block}{Buch, S.~153}
    Topic-Modelle erlauben, Texte nicht mehr nur \emph{sequenziell}
    zu lesen, sondern sie auch \emph{parallel} zu berücksichtigen.
  \end{block}

  \vfill
  Die klassische qualitative Forschung liest Texte einzeln,
  sequenziell, hermeneutisch. Topic-Modelle erlauben:
  \begin{itemize}
    \item parallele Sicht auf Hunderte oder Tausende von Texten,
    \vfill
    \item \emph{Navigation} im Korpus über die topics,
    \vfill
    \item \emph{vergleichende} Lektüre der besten Dokumente pro topic,
    \vfill
    \item Erkundung der \emph{semantischen Nachbarschaften} von Begriffen.
  \end{itemize}

  \vfill
  Das ist nicht der \emph{Ersatz} qualitativer Lektüre. Es ist ihr
  \emph{Vorlauf} -- eine Strukturierung, die das eigentliche
  Lesen vorbereitet.

\end{frame}

% ===============================================================
\begin{frame}{Konsequenzen für die Forschungspraxis}

  Wer mit Topic-Modellen arbeitet, ändert eine Reihe von Praktiken:

  \vfill
  \begin{itemize}
    \item \textbf{Korpusgröße:} Hunderte oder Tausende von Texten
          werden zugänglich. Was \enquote{relevante Literatur} ist,
          wird neu definiert.
    \vfill
    \item \textbf{Forschungsfrage:} sie muss explizit \emph{operationalisiert}
          werden -- in Form von Vorverarbeitungsentscheidungen und
          Parameterwahl.
    \vfill
    \item \textbf{Kooperation:} weil Parameter dokumentiert werden,
          können andere die Analyse \emph{reproduzieren}, \emph{kritisieren},
          \emph{weiterentwickeln}.
    \vfill
    \item \textbf{Wissenschaftliche Schreibweise:} man schreibt über
          \emph{Verfahren}, nicht nur über \emph{Befunde}.
          Methodenkapitel werden länger.
  \end{itemize}

\end{frame}

% ===============================================================
\begin{frame}{Was wir gelernt haben -- Bilanz der drei Sitzungen}

  \begin{enumerate}
    \item \textbf{Sitzung 1:} Topic-Modelle gehören zu den qualitativen
          Methoden -- sie teilen mit objektiver Hermeneutik und
          qualitativer Inhaltsanalyse drei Grundannahmen, aber sie
          unterscheiden sich in dem, was bei ihnen \emph{algorithmisch}
          ist.
    \vfill
    \item \textbf{Sitzung 2:} Mathematisch sind sie eine
          \emph{Matrixfaktorisierung} mit \emph{Mixed-Membership}.
          LDA und NMF sind zwei Wege zum gleichen Ziel. BERTopic
          und LLMs sind moderne Alternativen mit Vor- und Nachteilen.
          MTA bleibt aus guten Gründen bei LDA und NMF.
    \vfill
    \item \textbf{Heute, Sitzung 3:} Der Algorithmus ist nur die
          \emph{Mitte} einer Topic-Modell-Analyse. Vorverarbeitung
          und Interpretation machen den \emph{Großteil} der Arbeit
          aus -- und tragen die größte Verantwortung.
  \end{enumerate}

\end{frame}

% ===============================================================
\begin{frame}{Übergang zur Praxis}

  Ab der nächsten Sitzung beginnen wir die praktische Arbeit mit MTA.

  \vfill
  \textbf{Was Sie für die erste Praxis-Sitzung mitbringen sollten:}
  \begin{itemize}
    \item Ein kleines eigenes Korpus von Texten (5--20 Dateien)
          als \texttt{.txt} (UTF-8). \\
          \footnotesize Beispiele: Zeitungsartikel zu einem Thema, das
          Sie interessiert; ein paar Interviewtranskripte aus einem
          früheren Projekt; offene Antworten aus einer Umfrage.
          \normalsize
    \vfill
    \item Eine Frage, die Sie an diesen Texten stellen würden.
    \vfill
    \item MTA bereits auf Ihrem Laptop installiert (Installationsanleitung
          auf GitHub).
  \end{itemize}

  \vfill
  \begin{block}{Erinnerung}
    Die theoretischen Grundlagen, die wir besprochen haben, sind
    \emph{nicht} optional. Sie sind die Grundlage, auf der jede
    Topic-Modell-Analyse sinnvoll wird.
  \end{block}

\end{frame}

% ===============================================================
\begin{frame}[allowframebreaks]{Zitierte Literatur}

  \footnotesize

  \textsc{Grundlage:} \\[2pt]
  Papilloud, C., \& Hinneburg, A.\ (2018). \emph{Einführung in die
  qualitative Analyse von Texten mit Topic-Modellen}. Wiesbaden:
  Springer. \emph{Bes.\ Kap.\ 2.5 und Kap.\ 8.}

  \medskip
  \textsc{Vorverarbeitung von Textquellen:}
  \begin{itemize}
    \item Boyd-Graber, J., Mimno, D., \& Newman, D.\ (2014).
          Care and Feeding of Topic Models. In E.\ M.\ Airoldi et al.\
          (Hrsg.), \emph{Handbook of Mixed Membership Models}. CRC Press.
    \item Manning, C.\ D., Raghavan, P., \& Schütze, H.\ (2008).
          \emph{Introduction to Information Retrieval}. Cambridge UP.
          (TF-IDF, Kap.\ 6.2)
    \item Mimno, D.\ (2015). Optimizing Semantic Coherence in Topic
          Models. \emph{EMNLP}.
    \item Schofield, A., \& Mimno, D.\ (2016). Comparing Apples to
          Apple: The Effects of Stemmers on Topic Models. \emph{TACL}.
  \end{itemize}

  \medskip
  \textsc{Metadaten und erweiterte Topic-Modelle:}
  \begin{itemize}
    \item Blei, D.\ M., \& Lafferty, J.\ D.\ (2006). Dynamic Topic
          Models. \emph{ICML}.
    \item Roberts, M.\ E., et al.\ (2014). Structural Topic Models for
          Open-Ended Survey Responses. \emph{American Journal of
          Political Science}, 58(4), 1064--1082.
    \item Rosen-Zvi, M., Griffiths, T., Steyvers, M., \& Smyth, P.\
          (2004). The Author-Topic Model for Authors and Documents.
          \emph{UAI}.
    \item Mimno, D., \& McCallum, A.\ (2008). Topic Models Conditioned
          on Arbitrary Features with Dirichlet-Multinomial Regression.
          \emph{UAI}.
  \end{itemize}

  \medskip
  \textsc{Kritische Reflexion:}
  \begin{itemize}
    \item Blei, D.\ M.\ (2013). Topic Modeling and Digital Humanities.
          \emph{Journal of Digital Humanities}, 2(1).
    \item Gelman, A., Hennig, C., et al.\ (2013). Comment on
          \enquote{Hierarchical} Models in the Social Sciences.
    \item Mimno, D., \& Blei, D.\ (2011). Bayesian Checking for Topic
          Models. \emph{EMNLP}.
    \item Schmidt, B.\ M.\ (2012). Words Alone: Dismantling Topic
          Models in the Humanities. \emph{Journal of Digital
          Humanities}, 2(1).
  \end{itemize}

\end{frame}

% ===============================================================
\begin{frame}[plain]
  \begin{center}
    \vspace*{1.5em}
    {\Large\bfseries Ende des theoretischen Teils.}
    \\[2em]
    Ab der nächsten Sitzung: Praxis mit MTA.
    \\[3em]
    \footnotesize\color{black!60}
    christian.papilloud@soziologie.uni-halle.de
  \end{center}
\end{frame}

\end{document}
